% X_0^6(385)/<w_5,w_6,w_11,w_14> is hyperelliptic: a genus-three double cover of the genus-one star 2310q1,
% whose branch places pair exactly.  The method is that of paper/hyperelliptic-15-38.tex: the dual test gives the
% pairing modulo torsion, the component group at the Cerednik-Drinfeld prime 2 removes the torsion, and the Prym of
% the cover, simple over Q, shows the cover is the untwisted one.
% Tests: tests/Hyperelliptic6_385.m and tests/_offline/Hyperelliptic6_385.m.
\documentclass[11pt,reqno]{amsart}
\usepackage[margin=1.15in]{geometry}
\usepackage{amsmath,amssymb,amsthm,mathtools}
\usepackage[colorlinks=true,linkcolor=blue,citecolor=blue,urlcolor=blue]{hyperref}
\theoremstyle{plain}
\newtheorem{theorem}{Theorem}[section]
\newtheorem{lemma}[theorem]{Lemma}
\theoremstyle{definition}
\newtheorem{remark}[theorem]{Remark}
\DeclareMathOperator{\Jac}{Jac}
\DeclareMathOperator{\Pic}{Pic}
\newcommand{\Q}{\mathbb{Q}}
\newcommand{\Z}{\mathbb{Z}}
\begin{document}
\title{$X_0^{6}(385)/\langle w_5,w_6,w_{11},w_{14}\rangle$ is hyperelliptic}
\author{}
\date{\today}
\maketitle

\section{Statement}

Let $X=X_0^6(385)$, $W$ its Atkin--Lehner group of order $32$, $H=\langle w_5,w_6,w_{11},w_{14}\rangle$ of index
two, $Y=X/H$ of genus three and $E=X/W$ of genus one. The double cover $Y\to E$ is branched at the fixed points of
the involutions in $W\setminus H$: the rational CM points $P_{-40}$, $P_{-1540}$ and the place $Z_{-9240}$ of
degree two, the two CM points of discriminant $-9240$ on $E$.

\begin{theorem}\label{thm:main}
$Y$ is hyperelliptic: $P_{-40}+P_{-1540}\sim Z_{-9240}$ on $E$, and $Y$ is the double cover $E(\sqrt{cx})$ of a
function $x$ with that divisor.
\end{theorem}

The criterion and the method are those of \texttt{paper/hyperelliptic-15-38.tex}, Sections 1--4, to which we
refer for the proofs: a bielliptic curve of genus three is hyperelliptic if and only if its four branch points
split into two pairs with equal sums in the group $E$ and the cover is the untwisted one. The only Galois-stable
pairing here is $\{P_{-40},P_{-1540}\}\,|\,Z_{-9240}$, and we write $\Delta=P_{-40}+P_{-1540}-Z_{-9240}$.

\section{Proof}

\emph{The curve $E$.} Its traces of Frobenius agree with the class 2310q of Cremona's tables at every good
prime up to the Sturm bound $1152$ of level $2310$. The component groups computed below have orders $20$ at $2$
and $4$ at $3$, which in that class singles out 2310q1, $y^2+xy=x^3-991x-8359$, with
$E(\Q)_{\mathrm{tors}}\cong\Z/2$ generated by $T=(-26,13)$. At $2$ the reduction is split multiplicative of type
$\mathrm{I}_{20}$ and $T$ reduces to the singular point, so $T\notin E^0(\Q_2)$ and $E(\Q)_{\mathrm{tors}}$
injects into $\Phi_2\cong\Z/20$ (\texttt{paper/hyperelliptic-15-38.tex}, Lemma 3.1). At $3$ the torsion lies on
the identity component, so $\Phi_3\cong\Z/4$ gives only a consistency check.

\emph{$\Delta$ is torsion.} The dual test of \texttt{paper/hyperelliptic-15-38.tex}, Proposition 2.1, at
$M=4\cdot1155=4620$ (\texttt{tests/\_offline/Hyperelliptic6\_385.m}): a principal part with divisor $\Delta$ in the
dictionary of Guo--Yang pairs to zero with the $a_p(E)$-eigenspace of the cusp forms of weight $\tfrac32$
vanishing at the empty indices, and the other pairings are refused.

\emph{$\Delta$ is zero.} At $p=2$ the Cerednik--Drinfeld fibre of $X$ has components indexed by two copies of
the $96$ left ideal classes of the Eichler order of level $385$ in the definite quaternion algebra of
discriminant $3$ and singular points indexed by the $288$ classes of the level-$770$ order; the quotient by $W$
is a graph with $11$ vertices whose critical group is $\Z/20$. All three branch discriminants are even, so $2$
ramifies in each CM field and every branch point reduces to a singular point (Molina, Theorem 5.3); each lies on a
singular point whose branches are swapped by an involution of $W$ divisible by $2$, hence on a smooth point of a
merged component of $E$. The two points of $Z_{-9240}$ reduce to the two components carrying $P_{-40}$ and
$P_{-1540}$, one each, so $\Delta$ specialises to zero in $\Phi_2$, and $\Delta=0$ by the injectivity of the
torsion. At $p=3$, where $P_{-40}$ and $P_{-1540}$ are inert and $Z_{-9240}$ ramified, the class is zero as well.
Molina's theorems are stated for every prime dividing $D$, $p=2$ included.

\emph{The cover is untwisted.} The $L$-polynomial of the Prym of $Y\to E$ at $p=19$, from the point counts of
$Y$ over $\mathbb{F}_{19^k}$, $k\le3$, is $361x^4-10x^2+1$, irreducible over $\Q$, so the Prym is simple over
$\Q$. The twist of the hyperelliptic cover by the rational $2$-torsion class has a Prym isogenous to a product of
two elliptic curves over $\Q$ (the quotients $u'^2=u\,q_1(u)$, $u'^2=u\,q_2(u)$ for the factorisation
$q=q_1q_2$ over $\Q$ of the quartic model of $E$ in the coordinate $u=x$, which exists because
$\Pic^0(E)(\Q)[2]\neq0$), and a constant twist does not change that. So $Y$ is the untwisted cover and is
hyperelliptic. \qed

\end{document}
